For $n$ eligible, period-compatible counter sites with observations $y_j$ and
present-day model values $\widehat y_j$, define residual
$e_j=\widehat y_j-y_j$ and
\begin{align}
\operatorname{Bias}
 &=\frac1n\sum_{j=1}^{n}e_j,
\label{eq:validation-bias}\\
\operatorname{MAE}
 &=\frac1n\sum_{j=1}^{n}|e_j|,
\label{eq:validation-mae}\\
\operatorname{RMSE}
 &=\sqrt{\frac1n\sum_{j=1}^{n}e_j^2}.
\label{eq:validation-rmse}
\end{align}
With $\bar y$ and $\bar{\widehat y}$ denoting sample means, Pearson association
and the observed-on-modelled calibration slope are
\begin{align}
r_P
 &=\frac{\sum_j(y_j-\bar y)(\widehat y_j-\bar{\widehat y})}
 {\sqrt{\sum_j(y_j-\bar y)^2\sum_j(\widehat y_j-\bar{\widehat y})^2}},
\label{eq:validation-pearson}\\
b
 &=\frac{\sum_j(\widehat y_j-\bar{\widehat y})(y_j-\bar y)}
 {\sum_j(\widehat y_j-\bar{\widehat y})^2},
\qquad a=\bar y-b\bar{\widehat y}.
\label{eq:validation-calibration}
\end{align}
Spearman association $r_S$ applies the Pearson formula to average ranks. A
mean-observation null comparison is
\begin{equation}
R^2_{\mathrm{null}}
=1-\frac{\sum_j(y_j-\widehat y_j)^2}
        {\sum_j(y_j-\bar y)^2}.
\label{eq:validation-null-r2}
\end{equation}
A negative value means worse squared error than the observed-mean null. All
metrics are reported with site, temporal, direction, and screenline coverage;
spatially blocked holdouts are required when counters inform calibration.
